\documentclass[a4paper,11pt]{article}

\usepackage{luatexja}
\usepackage{luatexja-fontspec}
\usepackage{amsmath}
\usepackage{unicode-math}

\title{位相変化効果の導入計算}
\author{Hayato Ito}
\date{\today}

\begin{document}

    \maketitle

    \section{概要}
    励起状態密度を計算する際に、レーザーの位相変化の計算を導入する。

    \section{計算方法}
１ステップあたりの$\rho_{gg}$, $\rho_{ge}$の時間発展を次の式で表す。
\begin{align}
    &\dot{\rho_{gg}} = \Gamma \rho_{ee} + \Omega_{(t)} Im(\rho_{ge})\\
    &\dot{\rho_{eg}} = -(\frac{\Gamma}{2} + i\Delta)\rho_{eg} + i\frac{\Omega_{(t)}}{2}(\rho_{gg}-\rho_{ee})\\
    &\dot{\rho_{ge}} = \dot{\rho_{eg}}^*
\end{align}
これを４次ルンゲクッタ法で計算し、各ステップで$\rho_{gg}$, $\rho_{ge}$を更新する。この時の$\Gamma$,$\Delta$は以下の値である
\begin{align}
    \Gamma &= 2\pi \gamma_{(rad/s)}\\
    \Delta &= 2\pi f * n_{(rad/s)}, f = 1(MHz), n=-150, -149, ..., 150
\end{align}
    また$\Omega_{(t)}$は時間発展の計算ステップにおいて以下の２通り計算した。
\begin{align}
    (A)&\Omega_{(t)} = \Omega_0 e^{ikx_{t}}\\
    (B)&\Omega_{(t)} = \Omega_{(t-dt)} e^{ikv_{t-dt}dt}\\
    &k = 2\pi / \lambda
\end{align}
    以上の計算を、全体のステップの後半($t_{max}/2 < t \le t_{max}$)で行うことで、$\rho_{ee}$の積分を計算する
\begin{align}
    \rho_{ee} &= 1 - \rho_{gg}\\
    \int_0^{t_{max}} \rho_{ee} dt &= \sum_{n=0}^{t_{max}/dt} Re(\rho_{ee(n)}) * dt
\end{align}
    以上の計算を、$n=-150, -149, ..., 150$の全てのデチューニングに対して行うことで、スペクトル図を得る。

    \section{各種パラメータ}
        以上の計算を行う際の各種パラメータは以下のように設定した。
\begin{align}
    dt &= 25(ns)\\
    N &= 320000\\
    t_{max} &= dt * N = 8(ms)\\
    \Omega_0 &= 1 (rad/s)\\
        \lambda &= 313(nm)\\
        k &= 2\pi / \lambda(rad/s)\\
        f &= 1(MHz)\\
        \Delta &= [-150\cdot2\pi , 150\cdot2\pi](M\cdot rad/s)
    \end{align}

    \section{結果}
    デチューニング$\Delta$に対して、$\int_0^{t_{max}} \rho_{ee} dt$の値をプロットしたスペクトル図をA,Bの両方の方法計算し画像データを保存した。\\
    Aの結果：spectroscopy_A.png\\
    Bの結果：spectroscopy_B.png\\
    また、全時間範囲での粒子の位置をプロットした図をx_vs_time.pngとして保存した。\\



\end{document}